% GENERATED FILE. Edit canonical lessons, then run npm run generate:books.
% canonical-source-hash: fnv1a64:3bda08b88668578b
% canonical-lessons: GE-C24-denken, GE-C24-verstehen, GE-C24-lesen, GE-C24-schreiben

\chapter{Verbs of the Mind}
\label{ch:ge-mind-verbs}
\hlchaptermodality{24}

\begin{chapteropening}
\textbf{By the end of this chapter:} \emph{I can say I think, I understand, I read and I write in German, break a strong verb's vowel in the du and er forms, and say which of the four German borrowed rather than inherited.}

The chapter closes on schreiben, the one loan among four inherited verbs: the reader produces all four ich forms and all four er forms, names liest as the only vowel-breaker, and takes Manuskript apart into the Latin hand German never inherited plus the Latin verb it did. It reaches back past its own chapter to reassess Chapter 17's Hand, the Hand/manus gap, and Grimm's law.
\end{chapteropening}

\section[denken]{denken — the verb that was inside \emph{danke} all along}

\label{lesson:GE-C24-denken}



\begin{quote}
In Chapter 3 you learned that \textbf{danke} is built on a verb meaning \textquotedblleft{}to think.\textquotedblright{} Here is that verb, with its own forms at last.
\end{quote}

\begin{sounds}
\begin{itemize}
\raggedright
  \item \texttt{vowel-e-german} — a short, open \emph{e}.
  \item \texttt{ng-cluster} — \emph{n} before \emph{k} is the \emph{ng} of \emph{sing}: \textbf{denken} = \emph{DENG-ken}.
\end{itemize}
\end{sounds}

\subsection*{You'll want to know: denken}
\textbf{denken} means \textquotedblleft{}\textbf{to think},\textquotedblright{} with the same audible endings you drilled on \emph{wohnen}, \emph{machen} and \emph{lernen}:

\begin{itemize}
\raggedright
  \item \textbf{ich denke} — I think · \textbf{wir denken} — we think
  \item \textbf{du denkst} — you think · \textbf{ihr denkt} — you all think
  \item \textbf{er denkt} — he thinks · \textbf{sie denken} — they think
\end{itemize}

\begin{quote}
\textbf{Ich denke, also bin ich.} — \textquotedblleft{}I think, therefore I am.\textquotedblright{}
\end{quote}

It also carries Chapter 21 straight into a new sentence:

\begin{quote}
\textbf{Ich denke, es ist kalt.} — \textquotedblleft{}I think it is cold.\textquotedblright{}
\end{quote}

\begin{cousinweb}
\textbf{denken} is Proto-Germanic \emph{\textbf{*þankijaną}}. English \textbf{think} is that same inherited word — a loan in neither direction.

One consonant separates them, and you have been taught how to read it. With \textbf{Haupt} beside \emph{head} you met \textbf{Grimm's law}, which turned Proto-Indo-European \emph{*k} into Germanic \emph{h}. The same law turned PIE \emph{*t} into Germanic \emph{*þ}, the \emph{th} sound. English kept it and spells it \textbf{th}: \emph{think}. German later softened it to \textbf{d}.

Germanic \emph{*þankō} meant both \textquotedblleft{}a thought\textquotedblright{} and \textquotedblleft{}gratitude\textquotedblright{} — hence German \textbf{Dank} and \textbf{Gedanke}, and English \textbf{thank} beside \textbf{think}. To thank was to hold someone in thought. German also kept a sibling verb, \textbf{dünken}, \textquotedblleft{}to seem\textquotedblright{}; English had that one too and let it collapse into \emph{think}.
\end{cousinweb}

\begin{grammarlens}[title={Grammar Lens: also does not mean \textquotedblleft{}also\textquotedblright{}}]
In \emph{Ich denke, also bin ich}, the word \textbf{also} means \textquotedblleft{}\textbf{therefore}.\textquotedblright{} English \emph{also} is \emph{eall swā}, \textquotedblleft{}all so,\textquotedblright{} meaning \emph{likewise}; German \emph{also} is \emph{al} + \emph{so}, meaning \emph{thus}. Same bricks, different building. A German saying \textbf{also} is drawing a conclusion, not adding to a list.
\end{grammarlens}

\subsection*{Your turn}
Say these aloud:

\begin{itemize}
\raggedright
  \item \textquotedblleft{}ich denke, du denkst, er denkt\textquotedblright{}
  \item \textquotedblleft{}wir denken, ihr denkt, sie denken\textquotedblright{}
  \item \textquotedblleft{}Ich denke, also bin ich.\textquotedblright{}
  \item \textquotedblleft{}Ich denke, es ist kalt. Das Wetter ist gut.\textquotedblright{}
  \item the family — \textquotedblleft{}denken, danke, Dank, Gedanke\textquotedblright{}
  \item the shift — PIE \emph{t}, Germanic \emph{th}, German \emph{d}: think, denken
\end{itemize}

\begin{tcolorbox}[breakable,colback=teal!4,colframe=teal!35!black,title={Before you move on}]
Which chapter-3 word is built on \textbf{denken}? (\textbf{Danke}.) Is English \emph{think} a borrowing? (\textbf{No} — both inherited \emph{*þankijaną}.) Which law turned PIE \emph{*t} into the Germanic \emph{th} German made a \emph{d}? (\textbf{Grimm's law}, as in \emph{Haupt} and \emph{head}.) What does German \textbf{also} mean? (\textbf{Therefore}.)
\end{tcolorbox}
\section[verstehen]{verstehen — standing around a thing until you see it}

\label{lesson:GE-C24-verstehen}



\begin{quote}
You can say \textbf{ich denke}. Thinking happens in the \textbf{Kopf}; this is the verb for when it lands — and English built its own word from the same picture, separately.
\end{quote}

\begin{sounds}
\begin{itemize}
\raggedright
  \item \texttt{v-as-f} — German \textbf{v} is said \emph{f}, so \emph{ver-} is \emph{fair-}.
  \item \texttt{h-silent-lengthening} — the \textbf{h} is silent; it marks the vowel before it as long. \textbf{verstehen} = \emph{fair-SHTAY-en}.
\end{itemize}
\end{sounds}

\subsection*{You'll want to know: verstehen}
\textbf{verstehen} means \textquotedblleft{}\textbf{to understand},\textquotedblright{} with the ordinary endings:

\begin{itemize}
\raggedright
  \item \textbf{ich verstehe} · \textbf{wir verstehen}
  \item \textbf{du verstehst} · \textbf{ihr versteht}
  \item \textbf{er versteht} · \textbf{sie verstehen}
\end{itemize}

\begin{quote}
\textbf{Ich verstehe.} — said on its own, constantly. \textbf{Ich verstehe Deutsch.} — \textquotedblleft{}I understand German.\textquotedblright{}
\end{quote}

\begin{cousinweb}
Break it in two: \textbf{ver-} plus \textbf{stehen}, \textquotedblleft{}to stand.\textquotedblright{} \emph{Stehen} is English \textbf{stand}, the same inherited word, Proto-Germanic \emph{\textbf{*standaną}}. The prefix \textbf{ver-} is old and worn, and here it carries its ancient sense \textquotedblleft{}around, about.\textquotedblright{} So \emph{verstehen} is literally \textquotedblleft{}\textbf{to stand around}\textquotedblright{} a thing.

Now English. \textbf{Understand} is \emph{under} + \emph{stand} — and that \emph{under} is not \textquotedblleft{}underneath.\textquotedblright{} The usual explanation is that Old English \emph{under} also carried the sense \textquotedblleft{}among, between,\textquotedblright{} a cousin of Latin \emph{inter}; so \emph{understand} is \textquotedblleft{}to stand \textbf{among}\textquotedblright{} a thing. That reading is standard but not certain, and no honest account will claim otherwise.

Either way the important fact holds: \textbf{each language welded a prefix onto its own inherited verb \textquotedblleft{}to stand\textquotedblright{} and got a word for comprehension.} Neither borrowed the compound. Chapter 23 showed the same move with \textbf{grün}: Germanic and Latin each built \textquotedblleft{}green\textquotedblright{} from a verb for \textquotedblleft{}grow,\textquotedblright{} independently. This is that coincidence one family closer in.

The verb underneath is one of the family's biggest roots, PIE \emph{\textbf{*steh\textsubscript{2}-}}: German \textbf{der Verstand} (\textquotedblleft{}reason\textquotedblright{}), English \textbf{stand}, \textbf{stead}, \textbf{stall}, and — from Latin's \emph{stāre} — \textbf{state}, \textbf{station}, \textbf{stable}, \textbf{constant}.
\end{cousinweb}

\subsection*{Your turn}
Say these aloud:

\begin{itemize}
\raggedright
  \item \textquotedblleft{}ich verstehe, du verstehst, er versteht\textquotedblright{}
  \item \textquotedblleft{}Ich verstehe Deutsch.\textquotedblright{}
  \item both verbs — \textquotedblleft{}Ich denke. Ich verstehe.\textquotedblright{}
  \item the pieces — \textquotedblleft{}ver- + stehen\textquotedblright{} beside \textquotedblleft{}under + stand\textquotedblright{}
  \item the parallel — like \emph{grün} beside \emph{viridis}, built, not shared
  \item where it happens — \textquotedblleft{}der Kopf\textquotedblright{} — and what stands in it, \textquotedblleft{}der Verstand\textquotedblright{}
\end{itemize}

\begin{tcolorbox}[breakable,colback=teal!4,colframe=teal!35!black,title={Before you move on}]
What are the two pieces of \textbf{verstehen}? (\emph{\textbf{Ver-}} + \emph{\textbf{stehen}}.) Which English verb is \emph{stehen}? (\textbf{Stand}.) Did English borrow \emph{understand} from German? (\textbf{No} — each built its own compound.) How solid is the \textquotedblleft{}under = among\textquotedblright{} account? (\textbf{Standard, but not certain}.) Say \textquotedblleft{}I think\textquotedblright{} and \textquotedblleft{}I understand.\textquotedblright{} (\emph{\textbf{Ich denke. Ich verstehe.}})
\end{tcolorbox}
\section[lesen]{lesen — gathering, and the first verb that changes its vowel}

\label{lesson:GE-C24-lesen}



\begin{quote}
Every German verb so far has kept its vowel in every form. This one does not — and its older meaning is \textbf{gathering}.
\end{quote}

\begin{sounds}
\begin{itemize}
\raggedright
  \item \texttt{long-e} — a long \emph{e}, English \emph{ay} without the glide.
  \item \texttt{s-voiced} — a single \textbf{s} between vowels buzzes like \emph{z}: \emph{LAY-zen}.
\end{itemize}
\end{sounds}

\subsection*{You'll want to know: lesen}
\textbf{lesen} means \textquotedblleft{}\textbf{to read}.\textquotedblright{}

\begin{itemize}
\raggedright
  \item \textbf{ich lese} — I read · \textbf{wir lesen} — we read
  \item \textbf{du liest} — you read · \textbf{ihr lest} — you all read
  \item \textbf{er liest} — he reads · \textbf{sie lesen} — they read
\end{itemize}

\begin{quote}
\textbf{Ich lese Deutsch.} — \textquotedblleft{}I read German.\textquotedblright{}
\end{quote}

\begin{grammarlens}[title={Grammar Lens: the vowel that breaks in the middle}]
Four forms keep the \textbf{e} of \emph{lesen}. Two swap it for a long \textbf{ie}: \emph{du \textbf{liest}}, \emph{er \textbf{liest}}.

That is the mark of a \textbf{strong verb}: the vowel changes inside the word, in exactly two present-tense forms — \textbf{du} and \textbf{er/sie/es}. Nothing else moves; the endings stay the familiar \emph{-e / -st / -t}. It is a rule about a whole class of verbs, not about \emph{lesen}.
\end{grammarlens}

\begin{cousinweb}
\textbf{lesen} is Proto-Germanic \emph{\textbf{*lesaną}}, \textquotedblleft{}\textbf{to gather}.\textquotedblright{} That sense still lives in \textbf{die Weinlese}, the grape harvest — the \emph{wine-gathering}, built on the \textbf{Wein} of Chapter 11. Reading came second: you pick the letters up, one after another.

Now the trap. Latin \textbf{legere} also meant \textquotedblleft{}to gather,\textquotedblright{} also came to mean \textquotedblleft{}to read,\textquotedblright{} and stands behind \textbf{legible}, \textbf{lecture}, \textbf{lesson} and \textbf{collect}. It looks like one shared word. It probably is \textbf{not}: the sounds do not correspond as they should, and the standard account ties \emph{lesen} to a separate root, whose clearest cousin is Lithuanian \emph{lèsti}, \textquotedblleft{}to peck up grain.\textquotedblright{} Treat it as you treated \textbf{grün} beside \emph{viridis}.

English \textbf{read} is a third story: Old English \emph{rædan}, \textquotedblleft{}to advise,\textquotedblright{} which is German \textbf{raten} and \textbf{das Rätsel} (\textquotedblleft{}riddle\textquotedblright{}). \emph{Lesen} and \emph{read} are no relation.
\end{cousinweb}

\subsection*{Your turn}
Say these aloud:

\begin{itemize}
\raggedright
  \item \textquotedblleft{}ich lese, du liest, er liest\textquotedblright{} — hear the vowel break
  \item \textquotedblleft{}wir lesen, ihr lest, sie lesen\textquotedblright{} — and back
  \item \textquotedblleft{}Ich lese Deutsch. Ich verstehe Deutsch.\textquotedblright{}
  \item the old meaning — \textquotedblleft{}die Weinlese\textquotedblright{}
  \item the three so far — \textquotedblleft{}denken, verstehen, lesen\textquotedblright{}
  \item why \emph{lesen} is not \emph{legere} — the reason \emph{grün} is not \emph{viridis}
\end{itemize}

\begin{tcolorbox}[breakable,colback=teal!4,colframe=teal!35!black,title={Before you move on}]
Which two forms of a strong verb break their vowel, and what are they here? (\textbf{Du} and \textbf{er/sie/es} — \emph{\textbf{du liest, er liest}}.) What did \emph{lesen} mean first? (\textbf{To gather} — \emph{die Weinlese}.) Is it Latin \emph{legere}? (\textbf{Probably not} — the sounds do not correspond.) And English \emph{read}? (\textbf{No relation} — it meant \textquotedblleft{}to advise,\textquotedblright{} and is German \emph{raten}.)
\end{tcolorbox}
\section[schreiben]{schreiben — the one borrowed verb}

\label{lesson:GE-C24-schreiben}



\begin{quote}
Three verbs in, all three English's own blood relatives. The fourth breaks the run: \textbf{German borrowed it; English did not.}
\end{quote}

\begin{sounds}
\begin{itemize}
\raggedright
  \item \texttt{sch-sh} — \textbf{sch} is one sound, English \emph{sh}.
  \item \texttt{diphthong-ei} — \textbf{ei} is said \emph{eye}: \textbf{schreiben} = \emph{SHRY-ben}.
\end{itemize}
\end{sounds}

\subsection*{You'll want to know: schreiben}
\textbf{schreiben} means \textquotedblleft{}\textbf{to write}.\textquotedblright{}

\begin{itemize}
\raggedright
  \item \textbf{ich schreibe} — I write · \textbf{wir schreiben} — we write
  \item \textbf{du schreibst} — you write · \textbf{ihr schreibt} — you all write
  \item \textbf{er schreibt} — he writes · \textbf{sie schreiben} — they write
\end{itemize}

\begin{quote}
\textbf{Ich schreibe Deutsch.} · \textbf{Die Hand schreibt.}
\end{quote}

\begin{grammarlens}[title={Grammar Lens: strong, but not here}]
\emph{Lesen} taught you to watch \textbf{du} and \textbf{er}. Here: \emph{du schreibst}, \emph{er schreibt} — \textbf{the vowel does not move.} \emph{Schreiben} is strong (its past is \emph{schrieb}), but its present is smooth. So the rule is exact: a strong verb \emph{may} break its vowel there. Hold the pair: \emph{du liest}, but \emph{du schreibst}.
\end{grammarlens}

\begin{cousinweb}
\textbf{schreiben} is Latin \emph{\textbf{scrībere}}, and before that \textquotedblleft{}\textbf{to scratch}.\textquotedblright{} German took it in with the writing itself, and early: the loan passed through German sound changes — \emph{sc-} became \emph{sch-} — and joined the native strong verbs, its past becoming \emph{schrieb}. A late loan would have done neither.

English refused it. \textbf{Write} is Old English \emph{wrītan}, also \textquotedblleft{}\textbf{to scratch},\textquotedblright{} because runes were cut; German \textbf{reißen} and \textbf{ritzen} are its relatives. Two languages, two writing-words, both meaning \emph{scratch}. Latin's verb reached English anyway: \textbf{scribe}, \textbf{script}, \textbf{describe}, \textbf{transcript}.

And one closes a circle from Chapter 17. German \textbf{Hand} has \textbf{no Latin cousin} — Latin says \emph{manus} — yet \emph{manus} reached German as a loan: \textbf{Manuskript}. \emph{Manu-} is \emph{manus}, the hand; \emph{-skript} is \emph{scrīptum}, this verb. A \textbf{Manuskript} is written \emph{by hand}: the Latin hand-word German never inherited, bolted to the Latin writing-verb it did.
\end{cousinweb}

\subsection*{Your turn}
Say these aloud:

\begin{itemize}
\raggedright
  \item \textquotedblleft{}Die Hand schreibt. Ich schreibe Deutsch.\textquotedblright{}
  \item the four \textquotedblleft{}I\textquotedblright{} forms — \textquotedblleft{}ich denke, ich verstehe, ich lese, ich schreibe\textquotedblright{}
  \item the four \textquotedblleft{}he\textquotedblright{} forms — \textquotedblleft{}er denkt, er versteht, er liest, er schreibt\textquotedblright{} — only \emph{liest} broke
  \item three inherited, one borrowed — \textquotedblleft{}denken/think, verstehen/stand, lesen/gather\textquotedblright{}, then \textquotedblleft{}schreiben/scrībere\textquotedblright{}
  \item denken's \emph{d} — PIE \emph{t}, Germanic \emph{th}: Grimm's law
  \item \textquotedblleft{}Ich denke, also bin ich.\textquotedblright{} — and what \emph{also} does
  \item \emph{Manuskript} — \emph{manus} plus \emph{scrīptum}
\end{itemize}

\begin{tcolorbox}[breakable,colback=teal!4,colframe=teal!35!black,title={Before you move on}]
Say all four \textquotedblleft{}I\textquotedblright{} forms. (\emph{\textbf{Ich denke, ich verstehe, ich lese, ich schreibe}}.) Which is the loan? (\emph{\textbf{Schreiben}}, from \emph{scrībere}, \textquotedblleft{}to scratch\textquotedblright{} — as English \emph{write} also meant.) Which \textquotedblleft{}he\textquotedblright{} form breaks its vowel? (Only \emph{\textbf{liest}}.) Which half of \emph{Manuskript} did German never inherit? (\emph{\textbf{Manus}}, its own \textbf{Hand}.) And German \emph{also}? (\textbf{Therefore}.)
\end{tcolorbox}
